\begin{proposition}[Stationary and simultaneous address constructions]\label{ex:app-simultaneous-address}
Use the definitions and realization freedoms in
Definitions~\ref{ex:sec7-typed-graph}, \ref{ex:sec7-symbol},
\ref{ex:sec7-address} and Lemma~\ref{ex:sec7-realization},
and the coefficient rules of Proposition~\ref{ex:app-coefficient-rules}.
Fix $b\ge1$. Either take one symbol on a nonempty set of always
active terminals with a nonzero polynomial in their pair bits, or
take one or two symbols on terminals with nonzero linear activations,
including a terminal of every nonzero activation, such that each
symbol separately has a nonzero terminal polynomial at every
$x\ne0$. In the latter case the symbols may read different
compatible projections of the common alphabet and fixed stationary
external data. Require the active data to satisfy each prescribed
filter at every nonzero address.
Starting from fixed compatible typed terminal and stationary
data, one can adjoin finitely many primary auxiliaries with nonzero
linear activations and choose only incidences involving them so that
all the symbols equal $1$ at every $x\ne0$. All original data stay
fixed, no auxiliary is present at zero, and the required neutral
external incidences and splitting conditions are retained.
Only nonempty configurations and nonempty contraction outputs are
used. A nonzero iterated difference, hence a partitionability
certificate, suffices for the required nonzero terminal polynomial.
For the two final symbols the zero-total packets below supply these
terminal hypotheses using only separate projection surjectivity.
For the even stationary application, the original prime representatives
may be fixed at the outset. For a stabilized odd symbol, first fix the
abstract data and the finite budget, then choose an actual stage
carrying them, and only then choose prime representatives. Previously
fixed representatives may be retained when compatible with the full
type at that stage; arbitrary earlier primes are not asserted to
satisfy a later moving condition.
\end{proposition}
\begin{proof}\leavevmode
\paragraph{Total labels of the final seeds.}
Both final symbols have projected total label zero. The intrinsic
odd filter contains $1$: its fixed unit squareclasses, positive sign,
and moving condition $(D/\ell_0)=1$ are those of $1$. The even source
for $E^k$ is likewise normalized in the unit filter containing $1$.
Every nonzero coefficient, and hence every unit seed, lies in its
respective filter. Additional stationary tests at primes dividing
$k$ refine the common alphabet but are forgotten by the intrinsic
odd projection. They do not add conditions to that odd symbol.

Take repeated packets with total zero in the common elementary
label group, including sufficiently many copies of every label.
For either projection separately, keep the first $r-1$ anchors of
its unit seed and contract the nonempty remainder into the last
anchor. Since the total projected label is zero, the remainder has
the last seed label. This argument uses surjectivity of each
projection separately. It does not require that a fixed old odd
seed admit lifts whose total is zero in the common alphabet.
The whole packet already has that total, and contractions preserve
it. Thus the final twists satisfy the stationary tests as well as
both coefficient filters.
Concretely, choose a packet with an even number of copies of every
common label, and enough copies to provide the anchors of either
seed under its respective projection.  Its total label is zero.
For every nonzero activation $\lambda$, take a disjoint copy of
this whole packet and give all its vertices activation $\lambda$.
At every $x\ne0$ there are exactly $2^{b-1}$ active packets.
The active list is therefore nonempty, has total common label zero,
and contains enough anchors for each projected seed separately.
Its terminal polynomial for each symbol is nonzero by the
contraction certificate, regardless of the initially fixed bits.
It also contains a terminal of each required activation.  This
constructs the terminal input of the address lemma explicitly.
No simultaneous lifting of a prescribed pair of seed labels is
needed.  Every added network vertex is primary, so the total
common label remains zero at each address.

